% LTeX: language=en-GB enabled=true

\documentclass[../main.tex]{subfiles}

\begin{document}

\newenvironment{qanda}{\setlength{\parindent}{0pt}}{\smallskip}
\newcommand{\Q}{\bigskip\bfseries Q: }
\newcommand{\A}{\par\textbf{A:} \normalfont}

\subsection{Evaluation Against the Requirements}

The original specification, from the Analysis, is copied again below:

\begin{enumerate}
    \item The program should allow the user to create file backups on a server, run by the client themselves, 
    \begin{enumerate}
        \item The software should be able to run on both the client and server, with appropriate functionality for each situation,
        \item They should be able to work together, with a server being able to store the user’s backed up files, 
    \end{enumerate}
    \item The client program should run on modern Linux systems, usable on a laptop or desktop computer, 
    \begin{enumerate}
        \item An installable binary package should be created that is executable on amd64 Linux systems, 
        \item It should be very unlikely for the client to use more than 50MB of RAM for prolonged periods of time,
        \item It should be very unlikely for the client to use more than 5\% of CPU time for prolonged periods of time while running on a modern laptop; my client's processor, the Intel i9-12950HX, can be used for a test case,
    \end{enumerate}
    \item It should be usable with an intuitive command line interface, not requiring interactive input, 
    \begin{enumerate}
        \item All required functions of the program should be able to be passed in through command line flags and parameters, 
        \item There can be an optional interactive mode, but this should not allow any functionality than the non-interactive mode does not, 
        \item There should be no graphical interface, 
    \end{enumerate}
    \item The server should be able to run on a low-end laptop 10-15 years old, 
    \begin{enumerate}
        \item The program should avoid reading and writing to the system drives as much as possible, since most such devices use HDDs, 
        \item It should be very unlikely for the server program to use more than 200MB of RAM for prolonged periods of time,
        \item It should be very unlikely for the server program to use more than 10\% of CPU time for prolonged periods of time while running on such a device; my client's desired server's processor, a 2\textsuperscript{nd} generation Intel i5 processor, can be used as a test case, 
    \end{enumerate}
    \item All file transfers should be secure, so a man-in-the-middle would not be able to intercept data,
    \begin{enumerate}
        \item Any data intercepted should be of a format that the interceptor would not be able to decipher, or gain any useful information from, 
        \item The methods used should have their correctness checked,
    \end{enumerate}
    \item Text files stored on the server should be compressed such that they use a reasonable amount of storage, 
    \begin{enumerate}
        \item They should be at most as large as (or very close to) the uncompressed files, i.e. the file size should never increase following compression by more than 10 bytes,
        \item The client should be able to decompress any file when needed by the user, storing any needed details about the compression within the file, 
    \end{enumerate}
    \item File transfers should be verifiably correct, so data is not corrupted during transmission,
    \begin{enumerate}
        \item Once a file has been transferred to the server, it should be compared with the original file on the client device, 
        \item This checking should be with a much smaller version of the file than that which was originally sent for most files – it should be up to 10kB in size, 
    \end{enumerate}
    \item Generally, the program should work robustly,
        \begin{enumerate}
            \item Problems with user input or command line arguments should be gracefully handled,
            \item The program should not crash while running,
            \item Errors should be reported to the client, server, or both, depending on which is appropriate in the situation.
        \end{enumerate}
\end{enumerate}

This table, below, details the extent to which the specification points have been achieved.

\begin{longtable}{|c|p{0.9\linewidth}|}
\hline
Spec & To what extent has it been achieved? \\
\hline
\endfirsthead

\hline
Spec & To what extent has it been achieved? \\
\hline
\endhead

\hline
\endfoot

\hline
\endlastfoot

1a & This has been fully achieved. The same program can be run on the client and server, specifying which it should act as through a flag. The client will make a connection to the server running at the IP address specified by the client. \\
\hline

1b & This has been fully achieved. The client can send files to the server, and the server stores them under the ServerReceived directory, maintaining the file structure of the client. For example, if the client sends the file \texttt{/home/shayan/test1.txt}, it will be stored at \texttt{ServerReceived/home/shayan/test1.txt} on the server. \\
\hline

2a & This has been fully achieved. The program comes with a self contained executable binary for amd64 Linux systems. \\
\hline

2b & This has been fully achieved. In all of my testing the client program never used more than 40MB of RAM, below the limit of 50MB set. This is partly due to my choice of algorithms, for example, opting to use the simpler Huffman Coding rather than Prediction by Partial Matching, which would use more RAM. \\
\hline

2c & This has been fully achieved. In my testing the client program only exceeded the 5\% limit very briefly, which is within the bounds of the specification point. For the vast majority of its runtime, the program used less CPU time than that. Do note that this specification point is limited to my client’s laptop and its processor - CPU usage will vary greatly by device. \\
\hline

3a & This has been fully achieved. The server program (running the program with the server -s flag) accepts no interactive input and requires the port to be passed in as a command line argument. The client program (running the program with the client -cr or -cs flag) has all of its functionality accessible through command line arguments. \\
\hline

3b & This has been fully achieved. When running the program with the client -cs or -cr flag, the server IP, server port, and file paths can be specified interactively. This allows the user to correct errors identified by the program. However, no core functionality is gained through the interactive mode so the program can be used in scripts without any problems. \\
\hline

3c & This has been fully achieved. There is no graphical interface and the program only works through a console, as intended. \\
\hline

4a & This has been fully achieved. Files are compressed before being sent, and are stored in compressed form on the server, only decompressed on the client device after being retrieved. This minimises how much data the program reads and writes from the server’s hard drive. \\
\hline

4b & This has been fully achieved. In all of my testing the client program never used more than 80MB of RAM, below the limit of 200MB set. \\
\hline

4c & This has been fully achieved. In my testing the client program never exceeded the 10\% limit. In fact, it never exceeded 2\% processor usage. Do note that this specification point is limited to my client’s server laptop and its processor. \\
\hline

5a & This has been fully achieved. First, an asymmetric RSA key exchange is used to create public and private keys to share a symmetric AES key between the client and server devices. From that point on, all data is encrypted symmetrically with that AES key. \\
\hline

5b & This has been fully achieved. A standard C\# implementation of AES is used, which is widely recognised as correct, and I have checked key components of my RSA implementation through testing. \\
\hline

6a & This has been fully achieved. Plaintext files are almost always significantly compressed, and all other files are kept the same size. There is also functionality to ensure that text files never increase in size unnecessarily. \\
\hline

6b & This has been fully achieved. Files are decompressed by the client once they are received. This has been shown in testing, where the files are identical after receipt to before being backed up. \\
\hline

7a & This has been fully achieved. A hash of the file is sent along with the file itself, and the recipient verifies integrity by comparing hashes. This is implemented for both sending and retrieving files. \\
\hline

7b & This has been fully achieved. The hash is 256 bits (32 bytes) in size, which is much smaller than the 10kB limit. \\
\hline

8a & This has been fully achieved. If any user input is invalid, the user is informed of this, and in some cases is given the opportunity to rectify it. \\
\hline

8b & This has been fully achieved. In all of my testing, the program has never crashed. Errors are handled gracefully through exception handling and defensive programming. \\
\hline

8c & This has been fully achieved. The program never fails silently; there is always some form of error message displayed. \\
\hline
\end{longtable}

As justified in the table, all of my specification points have been fully achieved.

\subsection{Client Feedback}

To gather some feedback from my client, I conducted a final interview with them, to get their opinion on the final product made. 

\subsubsection{Final Client Interview}
\begin{qanda}
\Q Overall, what do you think of the backup utility that I’ve created? 

\A I think it’s broadly excellent. It does what I want it to do - I can use my old laptop as a server to store backups of files on my main laptop. I love that all the file transfers are encrypted, and that the files are sent compressed. It has a very intuitive command line interface, and so ties in really neatly with my workflows, which are normally based in the terminal. Having the single executable binary means that I can very easily install it system-wide, and I am really pleased that all of its functions are available through command line environments, since I want to incorporate this into a few of my system administration scripts. 

\Q Thank you so much! Are there any areas where you see shortfalls? 

\A There are some small improvements that I would suggest. Do note that these are mostly matters of personal preference. Firstly, when the server is running, if it is listening on the broader private network (i.e. not on the loopback interface) it reports its IP address as 0.0.0.0. While this is not incorrect, as it signifies that the server is accessible on both its private IP address on the private network, and from the wider internet, since I will only use the server from a private network, I would rather that it displays the private IP address of the laptop on the network. Currently, if I forget the server’s private IP address, I would need to run separate commands to determine that. Another point is that there is no indication on the server when a file has been successfully received. Though I can see this on the client, since the program terminates once all of the files have been sent, it would be nice to have some indication of this on the server too, for my own sanity. Finally, while I really like that there are no graphics, currently command line arguments are just specified sequentially. I think it could be most intuitive if they followed flags. For example, instead of just needing to enter the server IP address and port after the -cs [client send] flag, perhaps something like this: \textit{-cs -\--serverip [server IP] -\--serverport [server port]} could be more intuitive. [The client wrote the italicised text on paper during the interview.] 

\Q Great, that’s really good to know! Besides from opinions of preference, are there any features you think should have been implemented? 

\A Well, you implemented everything I told you to, and I am very happy with the outcome, particularly given the deadline I set you, which you have even exceeded! 

\Q Thanks! But if you were to request the software again, and you no longer needed to worry about your deadline, are there any additional features you would have asked for? 

\A Yes, there are actually a few points. I think it could have been helpful to have an option to back up all the files in a given directory at once. For example, I would like to be able to specify /home/sam/ and have the program backup my whole home directory. This isn’t a big deal though, as since it can be used as part of a script, I can easily create a script to do this for me. Another useful feature could be the ability to store multiple versions of a file on the server, so I can see the different versions of my files over time. Following on from this, when backing up a file, if it hasn’t changed since the last backup, there could be some functionality to just keep the old copy, reducing runtime, network load, and the strain on my old hard drives. This is also useful as a standalone point, even if old versions can’t be kept - if a file hasn’t changed since the last backup, it shouldn’t be re-sent ideally. Finally, while I really like the encryption implemented to prevent malicious parties being able to view the data I am sending, it could also have been useful to have some method of verifying the authenticity of communications between the client and the server, to ensure that a malicious actor can’t retrieve my backups! 

\Q Thank you for that - that’s really useful. To wrap up, could you rate the overall backup utility out of 10 based on its ease of use, functionality, and overall fit to your requirements. 

\A Of course! I would give it a 9/10 for ease of use. I love the purely command line based interface, that it can run with just one executable binary, and the fact that all needed options can be passed in through the command line. The only point of improvement here is that I would have preferred the command line options to be passed in through flags, as I pointed out before. For functionality, I would give it an 8/10. It contains absolutely everything I needed, but there is of course more that could be added, as I recently mentioned, in the grand scheme of things. I think it is a 10/10 for the fit to my requirements. It achieved every point I indicated at the start of the project, and all in the time period I had designated. I am very happy with the outcome overall!

\end{qanda}

\subsubsection{Evaluating the Interview}

Overall, I am very pleased that my client was impressed with the software produced. He seemed to really like the general functioning of the backup utility, and made many positive comments about points that he had brought up in the initial client interview being met. I was also very pleased to hear that he considers the program to have fully met its specification.

However, there were also constructive points mentioned, about potential improvements that could be made. These fall under two broad categories. Firstly, there are improvements based on what they would prefer. These points were that they would like:
\begin{enumerate}
    \item The server's private IP address displayed when it is listening on the private network,
    \item The server to give a printed indication whenever a file is successfully received, and
    \item The program to accept command line arguments with named flags, rather than just sequentially.
\end{enumerate}
Secondly, there are improvements that include adding new functionality. These suggestions were:
\begin{enumerate}
    \item To enable the utility to store backups of all the files under one directory natively, by just specifying the parent directory,
    \item To allow the user to store multiple versions of a file as backups, rather than just the most recent copy,
    \item To only send a file to the server if it has changed since the previous backup, rather than sending it to overwrite an identical file in some situations, and
    \item To incorporate some sort of verification that the client or server is who they claim to be, so malicious actors cannot impersonate them.
\end{enumerate}
The first group of improvements can be implemented really simply, by just making minor tweaks to the code, since the overall functionality is not changed, just the way it interfaces with the user. However, more substantial additions are needed for the functionality suggested in the second group. These will be explored in depth in the next section.

\subsection{Improving the Outcome}

As identified by my client, and described in the previous section, there are four points of functionality that could be added to improve the outcome if the problem was revisited.

Firstly, the backup utility should be able to store backups of all of the files under one directory natively. This should be fairly straightforward to implement, using the built-in C\# \texttt{Directory.GetFiles()} method to get the names of all of the files under the directory specified by the user. At that point, each file should be treated as it would be had the client entered it separately. There would be a need to differentiate between when the user is specifying a directory, under which all files should be backed up, or just a file itself to back up. This distinction could be made by requiring directories to be entered with a leading \texttt{/}. For example, specifying \texttt{/home/shayan/} would back up all of the files under the \texttt{/home/shayan/} directory, while specifying \texttt{/home/shayan/test1.txt} would only back up the \texttt{test1.txt} file.

The second point mentioned was allowing the user to store multiple versions of a file as backups. This is considerably more difficult to implement, as it would involve changing the way files are stored on the server. Currently, files are stored according to their name and location on the client device. One option for implementing this could be as follows. Files could be stored as directories containing various versions of that file (by timestamp) on the server. This would mean that, for example, the file \texttt{/home/shayan/test1.txt} would be stored on the server at \texttt{ServerReceived/home/shayan/test1.txt/1774347692}, where \textit{1774347692} is the UNIX timestamp when the backup is made. In this way, when retrieving a file, the server can access every stored the versions of a file, and send them to the user. Perhaps first the server can send a list of available timestamps to the client device (with the client translating it to a human-readable date format), and then the user can choose which backup timestamp they would like to have retrieved, sending that timestamp to the server, with the server then finally returning the file. Perhaps there could also be a command line flag to specify that the $n$\textsuperscript{th} most recent backup should be returned, to prevent a reliance on the interactive mode, and allow the program to continue being used in scripts. 

This makes it more important to have a way of deleting backups from the server. There could be specific types of messages that the client can send to the server to indicate that it should delete a given backup (either the $n$\textsuperscript{th} backup, the oldest backup, or the backup with a certain timestamp for instance). Additionally, the server program could be configured with command line flags to only store the most recent $n$ backups of any given file, with it deleting the oldest backup once a new backup has been received that brings the total number of backups for that file to $n+1$. These additions would help manage the disk space is used on the server.

The third point mentioned by my client was that a file backup should only be sent from the client to the server if there is not already an identical file stored on the server. For example, if \texttt{ServerReceived/home/shayan/test1.txt} exists on the server, and the user requests to back up \texttt{/home/shayan/test1.txt}, and the contents of the two files are identical, then the file should not be needlessly sent to the server. However, there needs to be a way to check this. The already implemented hashing algorithm could be used for this. The client would send the name of the file to back up, as usual, but then the server could send a hash of the most recent backup of the file, if one exists (otherwise it could send a series of 256 zero bits for example, or some other form of message indicating the file is not present). Then, the client would compute a hash of the file it wants to back up, and compare it to the hash it received from the server, and only send the actual file to the server if the hashes are not identical. If they are, it could send some form of message indicating that, so the server does not expect a file. This could take the form of a single bit sent from the client to the server. For instance a 0 could indicate that the files are identical and so the server should expect nothing more, and a 1 could indicate that the file to store will follow. 

The final point mentioned was verifying that the devices are who they claim to be during transmissions. For this, I would need to implement digital signatures and signed certificates. This would require me to slightly change the way my program is structured. Currently, whenever a client device connects to the server, the server generates a new pair of RSA keys, using that key pair to generate the AES key that will be used for the remainder of those communications. In order to be able to verify authenticity, I would need to use the same RSA algorithm to generate a new key pair. The private key will be kept secret, but the public key will be put inside the program code. I would then need to generate certificates for any client devices, and the server, that my client wishes to use, signing these with that private key that has been generated. The certificates can be as simple as containing the public key of the device the owns the certificate, and a hash of they key that is encrypted using the aforementioned private key. Then, upon each connection, the two devices would send each other their signed certificate. By decrypting the hash in the certificate with the known and trusted public key that has been put in the program source code, the receiver can compare this with the sender's public key in the certificate. If they match, then the receiver can be assured that the sender is genuine, since they will need to have had access to trusted private key at some point. One of the public keys can then be used to establish the shared AES key, with the rest of the communications carrying on as normal. Do note that for this to work, the server can no longer generate a new RSA key pair for each connection; instead, each device in the system will need to have a fixed, pre-generated key pair, so that their certificate can contain their public key. Otherwise, they would need a new certificate signed before each transaction!

\end{document}
